% =====================================================================
%  ДӨРӨВ. УДИРДЛАГА ЗОХИОН БАЙГУУЛАЛТ, ХҮНИЙ НӨӨЦИЙН ТӨЛӨВЛӨГӨӨ
% =====================================================================
\chapter{УДИРДЛАГА, ЗОХИОН БАЙГУУЛАЛТ, ХҮНИЙ НӨӨЦИЙН ТӨЛӨВЛӨГӨӨ}

\section{Удирдлага, зохион байгуулалтын бүтэц}

Төслийг төслийн менежерийн зарчмаар удирдана. Төслийн баг нь гурван
гишүүнээс бүрдэх бөгөөд гишүүн бүр үндсэн үүргийнхээ зэрэгцээ хэд хэдэн
чиглэлийг хавсран хариуцна:

\begin{itemize}
  \item \textbf{М.Мөнхдорж — Төслийн менежер / Системийн архитектор:}
        Төлөвлөгөө, хугацаа, төсвийн хяналт, технологийн шийдвэр,
        архитектурын зураглал, аюулгүй байдлын загвар.
  \item \textbf{Ү.Тамир — Backend хөгжүүлэгч:} REST API, өгөгдлийн сан,
        төлбөрийн интеграци, ажилтны хэсгийн сервер талын логик.
  \item \textbf{Б.Эрдэнэбаяр — Frontend хөгжүүлэгч:} Хэрэглэгчийн PWA,
        ажилтны dashboard, төлбөрийн интерфэйс, UI/UX дизайн.
\end{itemize}

\begin{figure}[h]
  \centering
  \begin{tikzpicture}[node distance=9mm and 12mm,
      box/.style={draw, rounded corners=1.5mm, fill=blue!8,
                  minimum width=3.4cm, minimum height=0.95cm,
                  align=center, font=\small},
      mgr/.style={box, fill=orange!20}]
    \node[mgr] (pm) {М.Мөнхдорж\\ТМ / Архитектор};
    \node[box, below left=12mm and -12mm of pm.south] (dev1) {Ү.Тамир\\Backend хөгжүүлэгч};
    \node[box, below right=12mm and -12mm of pm.south] (dev2) {Б.Эрдэнэбаяр\\Frontend хөгжүүлэгч};
    \draw[thick] (pm.south) -- (dev1.north);
    \draw[thick] (pm.south) -- (dev2.north);
    \draw[thick, dashed] (dev1.east) -- node[above, font=\scriptsize] {хамтран} (dev2.west);
  \end{tikzpicture}
  \caption{Төслийн багийн бүтцийн диаграмм}
  \label{fig:ch4-bag}
\end{figure}

Тестийн ажиллагаа (QA), маркетинг, борлуулалтын чиглэлийг багийн
гишүүд ээлжлэн хариуцаж, шаардлагатай тохиолдолд гадаад, дотоодын
худалдааны түншүүдээс (freelancer, дизайнерын үйлчилгээ)
дэмжлэг авна.

\section{Компанийн ажиллах хүчний тооцоо}

Хөгжүүлэлтийн 6 сарын хугацаанд багийн 3 гишүүн бүрэн цагаар
ажиллана. Нэвтрүүлгийн дараах үе шатанд багийн бүрэлдэхүүн
хадгалагдаж, ажиллагааны дэмжлэг, маркетингийн ажлыг давхар
гүйцэтгэнэ (Хүснэгт~\ref{tab:ch4-hunii-nuuts}).

\begin{table}[h]
  \centering
  \caption{Хүний нөөц, цалин хөлс (сар бүр)}
  \label{tab:ch4-hunii-nuuts}
  \small
  \begin{tabular}{@{}p{4.0cm}cccc@{}}
    \toprule
    \textbf{Албан тушаал} & \textbf{Орон тоо} & \textbf{Үндсэн цалин} &
    \textbf{Нэмэгдэл} & \textbf{Нийт цалин} \\
    & & (\textcurrency{}/сар) & хөлс (\textcurrency{}) & (\textcurrency{}/сар) \\
    \midrule
    Төслийн менежер / Архитектор & 1 & 1{,}600{,}000 & 160{,}000 & 1{,}760{,}000 \\
    Backend хөгжүүлэгч & 1 & 1{,}500{,}000 & 150{,}000 & 1{,}650{,}000 \\
    Frontend хөгжүүлэгч & 1 & 1{,}500{,}000 & 150{,}000 & 1{,}650{,}000 \\
    \midrule
    \textbf{Нийт} & \textbf{3} & \textbf{4{,}600{,}000} &
    \textbf{460{,}000} & \textbf{5{,}060{,}000} \\
    \bottomrule
  \end{tabular}
\end{table}

\section{Цалингийн зардлын төлөвлөгөө}

Хөгжүүлэлтийн 6 сарын хугацаанд цалингийн нийт зардал:

\begin{equation}
  5{,}060{,}000 \times 6 = 30{,}360{,}000 \text{ төгрөг}
\end{equation}

Ажилтны нийгмийн даатгалын шилжүүлгийг нийт сарын цалингийн
12.6\%-иар тооцвол сар бүр 637{,}560\textcurrency{} (5{,}060{,}000 × 12.6\%)
буюу 6 сарын хугацаанд 3{,}825{,}360\textcurrency{} болно. Ингэснээр
цалингийн зардлын нийт дүн (даатгалтай): 30{,}360{,}000 + 3{,}825{,}360 =
34{,}185{,}360\textcurrency{} байна.

\noindent\textbf{Тайлбар:} Дээрх тооцоо нь хөгжүүлэлтийн 6 сарын
хугацааг хамарсан бөгөөд нэвтрүүлгийн дараах үе шатанд багийн
бүрэлдэхүүн хадгалагдаж, зардлын хэмжээ нь орлогын өсөлттэй
уян хатан уялдана. Цалингийн хэмжээг багийн гишүүдийг оюутан
хөгжүүлэгчид гэж үзэж, төслийн эхний шатны түвшинд тогтоож,
ирээдүйд компанийн хувь эзэмших боломжийг хавсруулсан.Хөгжүүлэлтийн зардлыг төслийн санхүүгийн төлөвлөгөөний (Зургаадугаар
бүлэг) хөрөнгө оруулалтын хэсэгт бүрэн тусгасан.

\noindent\textbf{Ажиллах хүчний хөгжил:} Багийн гишүүдийн сургалт,
туршлагын солилцоог сар бүрийн 1 өдрийн дотоод сургалтаар, мөн NFC,
төлбөрийн интеграцийн чиглэлээр гадаад, дотоодын мэргэжилтнүүдийн оролцсон
вэб семинараар хангана.
